% MODELO DE SUBMISSÃO DE RESUMO PARA APRESENTACAO DE TRABALHO
% E DE PÔSTER NO ERMAC-PR 2025

% BASEADO NOS TEMPLATES OFICIAIS DA SBMAC PARA O CNMAC 2024.
% PARA MAIS INSTRUÇÕES SOBRE FIGURAS, TABELAS E EQUAÇÕES,
% BAIXE O ARQUIVO DISPONIVEL EM 
% https://www.overleaf.com/latex/templates/modelo-para-submissao-de-resumos-para-o-cnmac-2024/qtqxfdmgttng


\documentclass[a4paper, 11pt]{article}

\usepackage[brazil]{babel}      % para texto em Portugues
%\usepackage[english]{babel}    % para texto em Ingles

\usepackage[utf8]{inputenc}   % para acentuação em Português com o uso do Unicode, 

%% POR FAVOR, EVITE MUDAR A CONFIGURAÇÃO DESTE MODELO.
%% ADICIONE APENAS OS PACOTES ESSENCIAIS PARA SEU TRABALHO.

\usepackage[centertags]{amsmath}
\usepackage{amssymb,amsfonts,amsthm}
\usepackage{indentfirst}
\usepackage{array}
\usepackage{float}
\usepackage[usenames,dvipsnames,svgnames,table]{xcolor}
\usepackage{hyperref}
\usepackage{enumerate}
\usepackage{graphicx}
\usepackage{epstopdf}
\usepackage{graphics}
\usepackage{subfigure}
\usepackage[margin=1.5cm,footskip=1.5cm,includefoot]{geometry} 
\usepackage{listings}
\usepackage{setspace}
\usepackage{multirow}
\usepackage{multicol}
\usepackage{lipsum}

\setstretch{1.1}

\begin{document}
	%
	%
	%********************************************************
	% ************ CABEÇALHO - NAO MODIFICAR ****************	
	%********************************************************
	%
	%

	%
	%********************************************************
	%****************** TÍTULO DO TRABALHO ******************	
	%********************************************************	
	%
	%
	\begin{center}
	\textbf{Encontro de Matemática Aplicada e Computacional do Paraná | ERMAC-PR}
	\end{center}
	\vspace{-1.1em}
	\hrulefill\\
	\begin{center}
		\large{Passos para construir seu resumo e submeter no ERMAC-PR}
	\end{center}
	%
	%
	%********************************************************
	%*************** INFORMAÇÕES DOS AUTORES ****************	
	%********************************************************	
	%
	%
	% LEMBRE-SE DE SUBLINHAR O NOME DE QUEM VAI REALIZAR A
	% APRESENTAÇÃO
	%

	% PRIMEIRO AUTOR
	\noindent {\small 
	\underline{Juliana Verga Shirabayashi}\footnote{\texttt{juliana.verga@ufpr.br}},
	%
	% SEGUNDO AUTOR	
	Evandro Estevão Marquesone\footnote{\texttt{marquesone@utfpr.edu.br}},
	%
	% TERCEIRO AUTOR
	Glaucia Maria Bressan\footnote{\texttt{glauciabressan@utfpr.edu.br}}}\\
	%
	% ORIGEM DO PRIMEIRO AUTOR
	{\footnotesize
	$^{1}$ UFPR, Jandaia do Sul, PR\\ 
	%
	% ORIGEM DO SEGUNDO AUTOR
	$^{2,3}$ UTFPR, Cornélio Procópio, PR
	%
	}
	%
	%
	%********************************************************
	%**************** DESCRIÇÃO DO TRABALHO *****************	
	%********************************************************	
	%
	% A CRIAÇÃO DE SESSÕES É LIVRE, DESDE QUE RESPEITADO O
	% LIMITE DE DUAS (3) PÁGINAS, INCLUINDO AS REFERÊNCIAS.
	%

\begin{center}
\textbf{Resumo}
\end{center}

\indent Este é o formato padrão (apenas em \LaTeX) a ser usado exclusivamente nos resumos para trabalhos, que devem ter no máximo duas páginas, a serem submetidos no ERMAC-PR. Serão avaliados para sessões de pôsteres projetos desenvolvidos no nível de Iniciação Científica e, ou, trabalhos de pesquisa com resultados preliminares ou completos. Para Comunicação Oral são trabalhos no nível de Mestrado e, ou Doutorado com resultados preliminares ou completos. Os trabalhos submetidos que não estiverem de acordo com a formato padrão serão {\bf rejeitados} pelo Comitê Editorial do evento, sem análise do mérito científico.

 No cabeçalho deste documento, consta o(a) autor(a)-apresentador(a) sublinhado. Demais autores(as) vêm em seguida, de modo que todos(as) tenham indexados sua instituição de origem.
Os trabalhos podem ser submetidos em Português ou Inglês e serão apresentados na forma indicada pelo autor(a), pôster ou comunicação oral, desde que aprovada pela Comissão Editorial. Para que um trabalho aceito seja incluído na Programação do evento, é necessário que o(a) {\bf autor(a)-apresentador(a) tenha pago a taxa de inscrição até a data definida na página do ERMAC-PR}. Para submeter um trabalho na categoria de pôster, o(a) autor(a)-apresentador(a) deverá, no mínimo, estar cursando a graduação. Já para a categoria de comunicação oral, o(a) autor(a)-apresentador(a) deverá, no mínimo, estar cursando o mestrado ou doutorado. Cada autor(a)-apresentador(a) poderá submeter até dois resumos e o limite do resumo é de 3 páginas, para pôster e comunicação oral.  

{\bf As referências devem estar em ordem alfabética pelo sobrenome do primeiro autor.} Cada referência é produzida através do comando  \verb!\bibitem{nome-da-ref}! e é citada no texto via o comando \verb!\cite{nome-da-ref}!. A bibliografia ({\it Referências}), que deve figurar no final do artigo, é então gerada da seguinte forma (exige dupla compilação):

\begin{verbatim}
\begin{thebibliography}{00}
\bibitem{}
\end{thebibibliography}
\end{verbatim}

No caso de livros, deve-se seguir o padrão da referência \cite{Boldrini}, ou, para aqueles publicados dentro de uma séria, \cite{Gomes}. Se capítulo de livro, após o título da publicação, deve vir o título da série (quando aplicável), o número do capítulo e o volume, como na referência \cite{daSilva}; Se artigo, após o título da publicação ({\it em itálico}) deve vir o volume e as páginas correspondentes, seguidos do ano, conforme as referências \cite{Diniz2}. Trabalhos aceitos, mas não publicados, devem ser citados conforme mostrado na referência \cite{Cuminato}. Não havendo DOI, coloque-se (to appear). Trabalhos publicados em anais de eventos devem seguir o padrão da referência em \cite{Santos}. Dissertações, teses e similares devem seguir o padrão da referência \cite{Diniz1}. {\it Website} e outros \cite{SBMAC}.



As figuras devem ser inseridas de maneira que seu tamanho torne clara a sua leitura e compreensão, sem descaracterizar os objetivos propostos pelo texto. A legenda deve vir abaixo da mesma. \textbf{Quando não própria, deve-se indicar referência(s), fonte.} Para gerar uma figura como ilustrado pela Figura \ref{figura01}. 

\begin{figure}[h]
\centering
\includegraphics[width=.25\textwidth]{logo_sbmac.png}
\caption{ {\small Logo da SBMAC. Fonte: \cite{SBMAC}.}}
\label{figura01}
\end{figure}

Para a confecção das tabelas, deve-se usar o ambiente \texttt{table}, com a legenda acima da tabela e as entradas centralizadas nas colunas, como na Tabela \ref{tabela01}.

\begin{table}[ht]
\caption{ {\small Categorias dos trabalhos.}}
\centering
\begin{tabular}{ccc}
\hline
Categoria do trabalho  & Número máximo de páginas \\ 
\hline
Pôster          & 3  &     \\
Comunicação oral          & 3 \\
\hline
\end{tabular}\label{tabela01}
\end{table}


As equações são numeradas sequencialmente no texto, com a numeração automaticamente colocada à direita (favor não alterar) usando o comando \verb!\label{nome-da-equacao}! para identificá-las. A chamada \verb!\ref{nome-da-equacao}! faz referência à equação, no texto. Por exemplo, a equação (\ref{verhulst}) descreve o crescimento logístico,

%
\begin{equation}
\frac{dP(t)}{dt}= rP(t)\left(1-\dfrac{P(t)}{K}\right).
\label{verhulst}
\end{equation}




\noindent \textbf{Agradecimentos (opcional)}
\\
Seção reservada aos agradecimentos dos autores, caso pertinente.



	%
	%
	%********************************************************
	%********************* REFERÊNCIAS **********************	
	%********************************************************		
	%
	%
	% POR FAVOR, MANTENHA A BIBLIOGRAFIA NO PADRÃO RECOMENDADO.
	% COLOQUE APENAS AS PRINCIPAIS REFERÊNCIAS DO TRABALHO,
	% SEGUNDO A ORDEM ALFABÉTICA DO SOBRENOME DO PRIMEIRO AUTOR.
	%
	\begin{thebibliography}{00} \footnotesize
		\bibitem{Boldrini}
		Boldrini, J. L., Costa, S. I. R., Ribeiro, V. R., e Wetzler, H. G.. {\it Álgebra Linear e Aplicações}. Harper-Row, São Paulo, 1987.
		\bibitem{Cuminato}
		Cuminato, J. A. e Ruas, V., Unification of distance inequalities for linear variational problems, {\it Comp. Appl. Math.}, 2014. DOI: 10.1007/s40314-014-0163-6.

\bibitem{Diniz1}
Diniz, G. L. A mudança no habitat de populações de peixes: de rio a represa -- o modelo 
matemático, Dissertação de Mestrado, Unicamp, 1994.

\bibitem{Diniz2}
Diniz, G. L., Meyer, J. F. C. A. e Barros, L. C. Solução numérica para um problema de 
Cauchy Fuzzy que modela o decaimento radioativo, {\it TEMA},  23:63--72, 2001. DOI:10.1007/s40314-014-0163-6.

\bibitem{Gomes}
Gomes, L. T., Barros, L. C. e Bede, B. Fuzzy differential equation in various approaches. 
In {\it SpringerBriefs in Mathematics}. SBMAC- Springer, 2015. ISSN: 2191-8198.

\bibitem{Santos} Santos, I. L. D. e Silva, G. N. Uma classe de problemas de controle ótimo 
em escalas temporais, {\it Proceeding Series of the Brazilian Society of Computational and 
Applied Mathematics}, volume 1, 2013. DOI: 10.5540/03.2013.001.01.0177.

	\bibitem{daSilva} Silva, P. L. e Freire, I. L. On the group analysis of a modified Novikov equation, 
{\it Interdisciplinary Topics in Applied Mathematics, Modeling and Computational Science}, 
Springer Proceedings in Mathematics and Statistics,  volume 117, chapter 23, pages 161-166, 2015.

	\bibitem{SBMAC} SBMAC. Site oficial da Sociedade Brasileira de Matemática Aplicada e Computacional. Online. Acessado em 26/08/2024, \url{https://www.sbmac.org.br/}.

	\end{thebibliography}
\end{document}